\documentclass[11pt]{article}
\usepackage[margin=1in]{geometry}
\usepackage{amsmath,amssymb,amsthm,mathtools}
\usepackage{graphicx}
\usepackage{float}
\usepackage{microtype}
\usepackage[hidelinks]{hyperref}

\newtheorem{theorem}{Theorem}
\newtheorem{lemma}[theorem]{Lemma}
\newcommand{\R}{\mathbb R}
\newcommand{\C}{\mathbb C}
\newcommand{\N}{\mathbb N}
\newcommand{\Sph}{\mathbb S}

\title{A Forbidden Composite-Cosine Parameter\\Is Nevertheless Injective}
\author{Counterexample packet for arXiv:math/0501020}
\date{12 August 2026}

\begin{document}
\maketitle

\begin{center}
\fbox{\begin{minipage}{0.9\textwidth}
\textbf{Status: candidate full counterexample, likely valid; human review
requested.}
For \((n,m)=(3,2)\), the composite cosine transform with
\(\boldsymbol\lambda=(\alpha,0)\) is injective on
\(L^\flat(V_{3,2})\) exactly when \(\alpha\notin2\N_0\).  In particular,
\(T^{(1,0)}\) is injective even though both entries violate the condition
that Ournycheva and Rubin ask might be necessary when \(2m>n\).  Thus the
answer to their necessity question is no.
\end{minipage}}
\end{center}

\section{The source question}

For \(n>m\), Ournycheva and Rubin \cite{OR} define
\[
 (T^{\boldsymbol\lambda}f)(u)
 =\int_{V_{n,m}}f(v)\,(u'vv'u)^{\boldsymbol\lambda}\,dv,
 \qquad u\in V_{n,m},                                      \tag{1}
\]
on right-\(O(m)\)-invariant integrable functions.  If \(r\) is a positive
\(m\times m\) matrix and \(\Delta_j(r)\) is its leading principal
\(j\times j\) minor, its composite power is
\[
 r^{\boldsymbol\lambda}
 =\prod_{j=1}^m\left(\frac{\Delta_j(r)}{\Delta_{j-1}(r)}
                    \right)^{\lambda_j/2},
 \qquad \Delta_0=1.                                      \tag{2}
\]
In the convergence domain
\(\operatorname{Re}\lambda_j>j-m-1\), their Theorem~5.1 gives
injectivity under
\[
 \lambda_j+m-j\notin\{0,2,4,\ldots\},
 \qquad j=1,\ldots,m.                                    \tag{3}
\]
They prove (3) necessary for \(2m\le n\) and ask whether it is necessary
when \(2m>n\).

\begin{figure}[H]
\centering
\includegraphics[width=0.98\textwidth]{figures/open_problem_crop.png}
\caption{The source's question and surrounding conditions (PDF page 19).}
\end{figure}

\section{Idea of the proof}

Use the smallest over-rank Grassmannian, \(G_{3,2}\).  A plane is determined
by a unit normal \(w\), modulo sign.  On the parameter slice
\(\boldsymbol\lambda=(\alpha,0)\), the first principal minor of
\(u'vv'u\) is \(1-(w\cdot u_1)^2\), while the second exponent is zero.
The allegedly more complicated composite transform therefore becomes a
zonal convolution on \(Sph^2\).

The degree-\(2k\) multiplier of this convolution is an elementary Legendre
integral.  A beta integral and Chu--Vandermonde evaluate it as a product
whose only possible zero is \((-\alpha/2)_k\).  Hence every even harmonic
survives unless \(\alpha\) is a nonnegative even integer.  Choosing
\(\alpha=1\) makes all multipliers nonzero while making both exclusions in
(3) fail.

\section{The exact over-rank slice}

We use \((a)_k=a(a+1)\cdots(a+k-1)\) for the rising factorial.

\begin{lemma}[Legendre multiplier]
\label{lem:mult}
If \(\operatorname{Re}a>-1\) and \(k\in\N_0\), then
\[
 \int_{-1}^1(1-t^2)^aP_{2k}(t)\,dt
 =B\left(\frac12,a+1\right)
   \frac{(-a)_k(1/2)_k}{k!(a+3/2)_k}.                    \tag{4}
\]
\end{lemma}

\begin{proof}
Put \(s=t^2\) and use the terminating representation
\[
 P_{2k}(\sqrt{s})=(-1)^k\frac{(1/2)_k}{k!}
 {}_2F_1\left(-k,k+\frac12;\frac12;s\right).
\]
Termwise beta integration is legitimate because the series is a polynomial.
The beta ratio cancels the denominator parameter \(1/2\), leaving
\[
 (-1)^k\frac{(1/2)_k}{k!}B\left(\frac12,a+1\right)
 {}_2F_1\left(-k,k+\frac12;a+\frac32;1\right).
\]
Chu--Vandermonde turns the last factor into
\((a+1-k)_k/(a+3/2)_k=(-1)^k(-a)_k/(a+3/2)_k\), proving (4).
\end{proof}

\begin{theorem}[Complete criterion on a forbidden slice]
\label{thm:slice}
Let \((n,m)=(3,2)\), let \(\boldsymbol\lambda=(\alpha,0)\), and assume
\(\operatorname{Re}\alpha>-2\).  Then
\[
 T^{(\alpha,0)}:L^\flat(V_{3,2})\longrightarrow L^1(V_{3,2})
\]
is injective if and only if \(\alpha\notin\{0,2,4,\ldots\}\).
\end{theorem}

\begin{proof}
Identify a plane \(\eta\in G_{3,2}\) with a unit normal
\(w\in\Sph^2\), where \(w\) and \(-w\) represent the same plane.  Thus
\(L^\flat(V_{3,2})\) identifies, up to harmless normalization of measure,
with the even functions \(F\in L^1(\Sph^2)\).

Write \(u=(u_1,u_2)\in V_{3,2}\).  If the columns of \(v\) span
\(w^\perp\), then \(vv'=I-ww'\), and hence
\[
 \Delta_1(u'vv'u)=u_1'(I-ww')u_1=1-(w\cdot u_1)^2.       \tag{5}
\]
By (2), the second exponent vanishes when
\(\boldsymbol\lambda=(\alpha,0)\).  Consequently
\[
 (T^{(\alpha,0)}F)(u)
 =c\int_{\Sph^2}F(w)(1-(w\cdot u_1)^2)^{\alpha/2}\,d\sigma(w), \tag{6}
\]
where \(c>0\) depends only on the measure normalization.  The convergence
assumption is exactly \(\operatorname{Re}\alpha>-2\).

Set \(a=\alpha/2\).  By the Funk--Hecke formula, the operator in (6)
acts on every spherical harmonic of degree \(2k\) by a nonzero common
normalizing constant times the integral in (4).  In the half-plane
\(\operatorname{Re}a>-1\), the beta factor and all denominator factors in
(4) are finite and nonzero.  Thus the degree-\(2k\) multiplier vanishes
exactly when \((-a)_k=0\).

If \(a\notin\N_0\), no even multiplier vanishes.  If (6) is zero almost
everywhere, Fubini and Funk--Hecke show that every even spherical-harmonic
coefficient of \(F\) is zero.  Since spherical harmonics are uniformly
dense in continuous functions on \(\Sph^2\), and \(F\) is even, the finite
measure \(F\,d\sigma\) is zero; hence \(F=0\) almost everywhere.  The
transform is injective.

If \(a=q\in\N_0\), the kernel \((1-t^2)^q\) is a polynomial of degree
\(2q\).  Every nonzero spherical harmonic of even degree greater than
\(2q\) is therefore annihilated, so the transform is not injective.
\end{proof}

\section{Counterexample to the proposed necessity}

Take
\[
 n=3,\qquad m=2,\qquad \boldsymbol\lambda=(1,0).          \tag{7}
\]
This lies in the source's convergence domain because
\(1>-2\) and \(0>-1\).  But the two quantities excluded by (3) are
\[
 \lambda_1+m-1=2,
 \qquad
 \lambda_2+m-2=0,                                       \tag{8}
\]
so (3) fails in both coordinates.  Theorem~\ref{thm:slice}, with
\(\alpha=1\notin2\N_0\), nevertheless proves that \(T^{(1,0)}\) is
injective.  Therefore condition (3) is not necessary when \(2m>n\).

\section{Verification, scope, and novelty}

The argument answers the source's yes/no necessity question completely and
also determines injectivity on the whole slice
\((3,2,(\alpha,0))\).  It does not classify all multiparameters for every
pair \(2m>n\); after the negative answer, that sharper classification remains
a separate problem.

The accompanying script numerically integrates (4) at \(a=1/2\) for
\(k=0,\ldots,8\) and compares with the closed formula at 70-digit precision.
This checks transcription and signs but is not used as proof.  Expert review
should focus on the identification (5)--(6), the multiplier evaluation (4),
and the \(L^1\) completeness step.

On 12 August 2026, the four run indexes and bounded exact/current searches
for the source identifier, the exact quoted question, `composite cosine
transform', `\(2m>n\)', the authors, and polynomial-representation follow-up
found the source and its 2006 restatement, but no later explicit solution of
the multiparameter necessity question.  Novelty confidence is moderate-high.

\paragraph{Human-review recommendation.}
Review as a full counterexample.  The witness is explicit, lies strictly in
the source's convergence domain, violates both proposed exclusions, and its
injectivity follows from a complete elementary multiplier calculation.

\begin{thebibliography}{9}

\bibitem{OR}
E.~Ournycheva and B.~Rubin,
\emph{The Composite Cosine Transform on the Stiefel Manifold and Generalized
Zeta Integrals}, arXiv:math/0501020v2 (2005), Theorem~5.1 and the question on
PDF page~19.

\end{thebibliography}

\end{document}
